% ============================================================
% CHAPTER 7 — Internal Validation (REVISED)
% ============================================================

\chapter{Internal Validation}
\label{chap:internal_validation}

The preceding chapter defined five evaluation questions, of which EQ1
through EQ4 concern the correctness, scalability, and comparative
positioning of the framework against external references. This chapter
addresses EQ5, examining the internal validity of the methodology itself
before the comparative results are presented in
Chapter~\ref{chap:results}. The goal is to establish that the observed
behaviour of the framework is attributable to its design, namely the
combination of \ac{PPO} and \ac{DCR} shielding, rather than to either
component in isolation or to a particular experimental configuration.
Section~\ref{sec:shield_contribution} isolates the contribution of each
component through an ablation over three agent configurations, comparing
both the validity of the executions they produce and the quality of the
Pareto fronts they recover. Section~\ref{sec:rl_vs_safe_rl} examines the
training dynamics of the shielded and unshielded agents, showing that
the penalty signal alone does not converge to compliant behaviour.
Section~\ref{sec:policy_generalisation} assesses the robustness of the
learned policies across random seeds and discusses how generalisation is
appropriately evaluated in an online \ac{RL} setting. All experiments in
this chapter use the Loan Application benchmark and the training
protocol defined in Section~\ref{sec:training_config}.

% ------------------------------------------------------------
\section{Shield Contribution Analysis}
\label{sec:shield_contribution}

To isolate the contribution of each component, three agent
configurations are compared under an identical environment, reward
function, and interaction budget of 100\,000 timesteps with fixed seed
$s=1$, across the same six $(\alpha,\beta)$ weight configurations used
in the main experiments.

\begin{description}
    \item[Safe \ac{RL} (full framework).]
    \ac{PPO} with \ac{DCR} shielding, as described in
    Chapter~\ref{chap:method}. When the agent proposes an action whose
    event is not enabled under the current marking, the shield
    intercepts it before execution: the marking remains unchanged and a
    penalty of $-10$ is returned. This is the configuration used in all
    experiments of Chapter~\ref{chap:results}.

    \item[\ac{RL} only.]
    \ac{PPO} with the shield disabled. The penalty signal is preserved,
    a non-enabled action still returns $-10$, but the interception is
    removed: the \ac{DCR} engine applies the transition effects of the
    selected event unconditionally. An illegal action therefore mutates
    the marking, adding the event to the executed set, removing it from
    the pending set, and applying its include, exclude, and response
    relations, exactly as if it had been legally enabled. This
    configuration tests whether compliant behaviour can be learned from
    the reward signal alone.

    \item[Shielding only.]
    A non-learning agent that selects actions uniformly at random from
    the set of currently enabled events. This baseline quantifies the
    optimisation quality achievable from structural compliance alone,
    without any learning component.
\end{description}

The comparison admits two falsifiable predictions: if compliance were
learnable from the penalty signal alone, \ac{RL}~only should approach
the compliance of the full framework over training; if the shield
contributed only safety and nothing to optimisation, Shielding~only
should recover an equivalent Pareto front. Neither prediction holds, as
the results below show.

\paragraph{Structural validity of accepting episodes.}
An accepting episode can be invalid in two distinct senses, so two
metrics are needed. \emph{Executed-trace validity} asks whether any
transition actually applied to the marking was illegal at the time it
was executed; this is the metric tied to the compliance guarantee of
contribution~C1. \emph{Intervention-free} asks whether the episode
contains no blocked \emph{attempt} at all; under shielding, a blocked
attempt never touches the marking, so an episode can be intervention-rich
yet fully valid. The two are not interchangeable: blocked attempts leave
no trace in the executed behaviour, so only the first metric certifies
the process execution itself. Table~\ref{tab:ablation_validity} reports
both metrics for the three configurations.

\begin{table}[H]
\centering
\caption{Validity of accepting episodes per ablation configuration (final 20\% of training/evaluation). For Safe RL and RL only, the six $(\alpha,\beta)$ weight configurations are pooled, since each weight pair yields a distinct learned policy. Shielding only is a non-learning random-valid-action baseline, independent of $(\alpha,\beta)$, and is reported from a single evaluation run.}
\label{tab:ablation_validity}
{\renewcommand{\arraystretch}{1.3}
\setlength{\tabcolsep}{4pt}
\begin{tabular}{lrcc}
\toprule
\textbf{Configuration} & \textbf{Accepting episodes} &
\textbf{Executed-trace validity} & \textbf{Intervention-free} \\
\midrule
Safe \ac{RL}    & 17\,295 & 100\,\% & 93.8\,\% \\
\ac{RL} only    & 43\,986 & 0\,\%   & 0\,\%    \\
Shielding only  & 2\,900  & 100\,\% & 100\,\%  \\
\bottomrule
\end{tabular}}
\end{table}

The contrast is categorical. Under Safe \ac{RL} and Shielding only,
every accepting episode is a valid execution by construction, since the
shield prevents any illegal transition from reaching the marking. Under
\ac{RL} only, not one of the 43\,986 accepting episodes is structurally
valid: the acceptance condition (Equation~\ref{eq:accepting}) is
evaluated on the marking regardless of how that marking was produced,
so once illegal transitions are allowed to mutate the state, the
accepting label no longer certifies a valid process execution.

The intervention-free rate of the shielded agent, 93.8\,\%, measures a
different quantity, namely how completely the agent has internalised
the legal structure of the process by the end of training. The
remaining 6.2\,\% of accepting episodes contain at least one blocked
attempt, after which the agent proceeded to a valid accepting state.
These episodes are fully compliant at the level of executed behaviour;
the residual attempts are examined in
Section~\ref{sec:rl_vs_safe_rl}.

\paragraph{Pareto front comparison.}
Table~\ref{tab:ablation_pareto_summary} summarises the Pareto fronts
recovered by each configuration, and
Figure~\ref{fig:ablation_pareto_fronts} shows the dominance-filtered
fronts over the accepting episodes of the final 20\,\% of training.
Figure~\ref{fig:ablation_scatter} shows the full cost--duration
distribution of accepting episodes for the three configurations. The
\textbf{Safe \ac{RL}} row reproduces the same front already reported as
the headline result of Chapter~\ref{chap:results}
(Table~\ref{tab:loanapp_pareto_summary}); it is repeated here only as
the reference point against which the other two configurations are
compared.

\begin{table}[H]
\centering
\caption{Pareto fronts recovered per ablation configuration (final
20\,\% of training, all six weight configurations pooled). The Safe
\ac{RL} row is the same front reported in
Table~\ref{tab:loanapp_pareto_summary} of Chapter~\ref{chap:results}.}
\label{tab:ablation_pareto_summary}
{\renewcommand{\arraystretch}{1.2}
\setlength{\tabcolsep}{4pt}
\begin{tabular}{lrrcc}
\toprule
\textbf{Configuration} & \textbf{Accepting} & \textbf{Pareto points} &
\textbf{Cost range (a.u.)} & \textbf{Duration range (min)} \\
\midrule
Safe \ac{RL}   & 17\,295 & 26 & 1\,092--3\,120 & 408--1\,632 \\
\ac{RL} only   & 43\,986 & 2  & 14--40         & 7--28       \\
Shielding only & 2\,900  & 21 & 1\,092--3\,021 & 490.5--1\,632 \\
\bottomrule
\end{tabular}}
\end{table}


\begin{figure}[H]
    \centering
    % TODO: confirmar ruta
    \includegraphics[width=\linewidth]{Pictures/Resultado2_Pareto.png}
    \caption{Dominance-filtered Pareto fronts of the three ablation
    configurations on the Loan Application benchmark, over the
    accepting episodes of the final 20\,\% of training.}
    \label{fig:ablation_pareto_fronts}
\end{figure}

\begin{figure}[H]
    \centering
    % TODO: confirmar ruta
    \includegraphics[width=\linewidth]{Pictures/Result1_Scatter.png}
    \caption{Cost--duration distribution of all accepting episodes per
    ablation configuration (final 20\,\% of training). The shielded
    learning agent concentrates its episodes in bands near the
    non-dominated region, while the random shielded baseline disperses
    across the dominated region, and the unshielded agent populates a
    cost--duration range unreachable by any valid execution.}
    \label{fig:ablation_scatter}
\end{figure}

Taken at face value, the \ac{RL}-only front dominates the other two
numerically: both of its points lie an order of magnitude below the
cost and duration ranges recovered by the compliant configurations.
Inspection of the underlying episodes shows that this front is an
artefact of state corruption rather than an optimisation result. Both
points correspond to single-step episodes consisting of exactly one
illegal action, and both fire \textit{Cancel application} under the two
role assignments (Junior: cost 14, duration 28; Expert: cost 40,
duration 7). The same two points are recovered under every one of the
six weight configurations: regardless of $(\alpha,\beta)$, the
unshielded agent converges to this single exploit.

The mechanism is fully determined by the structure of the process
graph. In the initial marking of the Loan Application graph, the
pending set is empty and all twelve events are included.
\textit{Cancel application} is guarded by condition relations requiring
three other events to have been executed first, so it is not enabled at
step zero, and firing it as the first action is an illegal transition.
Crucially, \textit{Cancel application} has no outgoing response
relation: executing it adds nothing to the pending set, which therefore
remains empty. The acceptance condition is consequently satisfied after
this single step, and the episode terminates as accepting with the cost
and duration of one event, independently of the other eleven events in
the process. Under shielding this path is structurally unreachable,
since the shield blocks \textit{Cancel application} as a first action.
The two points of the \ac{RL}-only front are therefore two
role-conditioned instances of the same one-step exploit, describing a
process in which a loan application is cancelled before it exists. They
are excluded from the optimisation comparison on these grounds.

That the unshielded agent converges onto this exploit is consistent
with the reward structure: the $-10$ penalty is more than offset by the
terminal acceptance bonus, making the one-step shortcut the
highest-return behaviour available, so \ac{PPO} concentrates
probability mass on it. The exploit also explains the largest accepting
count of the three configurations, 43\,986: corrupted episodes are
extremely short, so many more fit within the same budget. This finding
illustrates, in an extreme form, that the \ac{DCR} acceptance condition
is a purely formal property of the marking; its implications for the
construct validity of the evaluation are discussed in
Chapter~\ref{chap:discussion}.
% TODO: ajustar la referencia anterior a la sección concreta de
% limitaciones (§ de construct validity / acceptance semantics).

With the \ac{RL}-only front excluded, the informative comparison is
between the full framework and the random shielded baseline, and it
isolates the contribution of the learning component. The two recover
near-identical fronts in the interior of the trade-off region: 21 of
the points found by the random baseline coincide with points on the
26-point front of the full framework. On this benchmark, then, a
substantial share of the front coverage comes from the structure of the
constrained trace space itself: restricted to valid executions, even
undirected sampling encounters most of the interior trade-offs.

However, because the Shielding-only baseline is stochastic, the
single-run result was repeated under four random seeds, using the same
evaluation budget as the original run. Table~\ref{tab:shielding_seed_variance}
reports the resulting seed variance, and Figure~\ref{fig:shielding_seed_variance} visualises the recovered fronts.

Across seeds, Shielding-only recovers $19.2 \pm 1.5$ Pareto points and
$2\,900.0 \pm 0.0$ accepting episodes. The fixed accepting count follows
from the fixed evaluation budget, while the recovered Pareto front
varies with random exploration. Figure~\ref{fig:shielding_seed_variance}
shows that the random shielded baseline repeatedly recovers the same
broad decreasing trade-off shape, but not the same front. The
minimum-cost endpoint is recovered in three of the four seeds, whereas
the minimum-duration endpoint varies more strongly, ranging from 408.0
to 530.0. Thus, Shielding-only can occasionally sample an extreme point
also recovered by Safe \ac{RL}, but this behaviour is not stable across
seeds.

\begin{table}[H]
\centering
\caption{Seed variance of the Shielding-only baseline on the Loan
Application benchmark. Each run uses the same random-valid-action
evaluation budget as the original Shielding-only run.}
\label{tab:shielding_seed_variance}
{\renewcommand{\arraystretch}{1.5}
\setlength{\tabcolsep}{6pt}
\begin{tabular}{rrrrc}
\toprule
\textbf{Seed} & \textbf{Accepting} & \textbf{Pareto points} &
\textbf{Cost range} & \textbf{Duration range} \\
\midrule
1 & 2\,900 & 21 & 1\,092.0--3\,021.2 & 490.5--1\,632.0 \\
2 & 2\,900 & 20 & 1\,092.0--4\,138.0 & 530.0--1\,632.0 \\
3 & 2\,900 & 18 & 1\,092.0--3\,246.0 & 430.5--1\,632.0 \\
4 & 2\,900 & 18 & 1\,148.7--3\,120.0 & 408.0--1\,669.5 \\
\midrule
Mean $\pm$ std & $2\,900.0 \pm 0.0$ & $19.2 \pm 1.5$ & -- & -- \\
\bottomrule
\end{tabular}}
\end{table}

This seed variance refines the interpretation of the ablation. The
shield alone is sufficient to guarantee structural validity and to
restrict exploration to the feasible process space, but it does not
provide a mechanism for reliably targeting a preferred region of the
trade-off. By contrast, the full framework recovers a larger and more
stable front, with $26.25 \pm 0.5$ Pareto points in the robustness analysis of Section~\ref{sec:policy_generalisation}
(Table~\ref{tab:seed_variance}), and its accepting episodes concentrate near the regions favoured by the scalarization weights. Learning therefore contributes
not merely additional points, but steerability: a random compliant agent
can discover valid trade-offs by chance, whereas a trained shielded
agent can be directed towards the operator's cost--duration preference.


\begin{figure}[H]
    \centering
    \includegraphics[width=0.75\linewidth]{Pictures/shielding_seed_variance_pareto_overlay.png}
    \caption{Pareto fronts recovered by the Shielding-only baseline
    across four random seeds on the Loan Application benchmark. The
    baseline samples uniformly among currently enabled actions and does
    not learn; variation across seeds therefore reflects stochastic
    coverage of the valid trace space under the same evaluation budget.}
    \label{fig:shielding_seed_variance}
\end{figure}




In summary, neither component suffices alone: without the shield the
agent optimises a corrupted process whose accepting states are not
valid executions, and without learning the shield guarantees validity
but leaves Pareto recovery stochastic, less stable, and not steerable.
Their combination provides validity by construction, directed coverage
of the valid front, and sample-efficient generation of valid traces, the
joint claim of contributions~C1 and~C2 that EQ5 was designed to test.
% ------------------------------------------------------------
\section{\ac{RL} versus Safe \ac{RL}: Training Dynamics}
\label{sec:rl_vs_safe_rl}

The previous section established what each configuration produces at the end of training. This section examines how the shielded and unshielded agents behave during training, addressing the question of whether the penalty signal alone can guide the model towards compliant behaviour. Figure~\ref{fig:illegal_rate_deciles} shows the average proportion of illegal actions per episode, pooled across the six weight configurations and binned into deciles of training progress. The shielding-only baseline is omitted here, as its illegal rate is zero by construction.

\begin{figure}[H]
    \centering
    % TODO: confirmar ruta
    \includegraphics[width=\linewidth]{Pictures/InternalVal_Result3_IllegalAction.png}
    \caption{Illegal action rate per training decile for the shielded and unshielded configurations, averaged across the six $(\alpha,\beta)$ weight configurations. For the shielded agent an illegal action is a blocked attempt; for the unshielded agent it is an executed illegal transition.}
    \label{fig:illegal_rate_deciles}
\end{figure}

The two curves differ significantly in both shape and final convergence. The shielded agent starts at 32.0\,\% in the first decile, rapidly drops below 4\,\% by the third decile, and stabilises at a residual 1.0\,\% by the end of training. In contrast, the unshielded agent starts at a much higher 82.4\,\%, decreases to roughly 50\,\% by the third decile, and enters a long plateau before finishing training with more than a third of its actions being illegal. While both agents learn to reduce violations over time, the unshielded agent remains orders of magnitude above the shielded baseline and shows no signs of full compliance within the given budget.

Two features of this comparison require explanation: the persistent plateau of the unshielded agent, and the large gap at the very start of training.

\paragraph{The plateau.}
Under shielding, an illegal attempt always yields a $-10$ penalty and blocks any environment progress, providing a consistent incentive to avoid them; the remaining 1\,\% error likely reflects standard exploration noise. The unshielded plateau, however, is not caused by the agent exploiting the process termination. A decile-level decomposition reveals that acceptance-inducing illegal actions (such as the \texttt{Event\_10} exploit described in Section~\ref{sec:shield_contribution}) actually decrease over time to below 1\,\%, representing only 6.9\,\% of total violations.

Instead, 91\,\% of late illegal actions are premature attempts to execute \textit{Design loan offer} and \textit{Approve loan offer} at the very beginning of almost every episode, without satisfying their DCR prerequisites. This suggests a persistent, sub-optimal opening habit. Because the flat $-10$ penalty does not scale with event consequences, it fails to provide a strong enough gradient to unlearn these expensive mistakes, leaving the agent trapped in a local minimum. Whether a cost-scaled penalty mechanism would effectively guide the agent out of this behavior is left for future work.

\paragraph{The initial gap.}
Both configurations share the same network architecture, hyperparameters, and random seed, meaning their initial policies are identical. Therefore, the wide gap in the first decile (82.4\,\% versus 32.0\,\%) is not an initialisation defect, but a result of environment dynamics.

First, the state distributions diverge immediately. In the unshielded setup, illegal actions apply invalid transitions that dynamically alter the DCR marking, shrinking the set of enabled events and making subsequent random choices more likely to be illegal as well. This compounding effect is absent under shielding, where the marking is structurally forced to stay on valid paths. Second, because the metric is an average over episodes, the extremely short length of corrupted episodes artificially inflates the per-decile violation rate.

\paragraph{Summary.}
To contextualise these findings within the core objectives of this thesis, Table~\ref{tab:ablation_compliance} consolidates the compliance metrics recorded during the final 20\,\% of training.

\begin{table}[H]
\centering
\caption{Compliance comparison between the shielded and unshielded configurations on the Loan Application benchmark (final 20\,\% of training).}
\label{tab:ablation_compliance}
{\renewcommand{\arraystretch}{1.5}
\setlength{\tabcolsep}{6pt}
\begin{tabular}{lcc}
\toprule
\textbf{Metric} & \textbf{\ac{RL} only} & \textbf{Safe \ac{RL}} \\
\midrule
Illegal action rate, final decile & 35.8\,\% & 1.0\,\% \\
Executed-trace validity of accepting episodes & 0\,\% & 100\,\% \\
Invalid transitions applied to the marking & present throughout & 0 \\
\bottomrule
\end{tabular}}
\end{table}

This structured comparison supports the conclusion central to contribution~C1: a simple penalty signal is insufficient to guarantee compliant behaviour. Compliance as a statistical tendency and compliance as a structural guarantee are not interchangeable; only the latter ensures that accepting traces represent genuinely valid executions of the business process.

% ------------------------------------------------------------
\section{Policy Generalisation}
\label{sec:policy_generalisation}

This section evaluates three robustness questions in turn: whether the
\ac{DCR} model is faithful to the real process (conformance replay),
whether the learned policies overfit the specific training
parameterisation (temporal holdout), and whether results depend on
optimisation stochasticity (random seeds).

\paragraph{Model faithfulness: conformance replay.}

The shield is synthesised from the \ac{DCR} model, so its guarantees
are only as meaningful as the model's fidelity to the behaviour of the
process it represents. To assess this fidelity, all 1\,000 traces of
the event log were replayed on the \ac{DCR} engine from the initial
marking, checking at every step whether the observed event was enabled
before applying its effects.

Of the replayed traces, 93.6\,\% contain no illegal transition, and
the two mainline variants of the process, comprising six and eight
events and accounting for 876 of the 1\,000 cases, replay without a
single violation. This non-conformance is unrelated to the role dimension: the enablement check that determines whether a transition is legal depends only on the \ac{DCR} marking (included, executed, and pending sets) and the condition, response, and exclusion relations of the graph, not on which
resource role executes an event. Roles affect only the cost and duration
of an executed transition, never its legality. The residual 6.4\,\% is
therefore a structural limitation of how the rework loop is modelled in
the \ac{DCR} graph, not an artefact of the role-conditioned action space.

The \ac{DCR} model is therefore faithful to the mainline behaviour of
the process, while capturing the rework behaviour only partially. This
delimits the scope of the compliance guarantee: the shield enforces
the modelled behaviour, and the fidelity of that model to the real
process is a separate, measurable property.

The replay also provides an empirical counterpart to the exploit
analysis of Section~\ref{sec:shield_contribution}. Across the 1\,000
observed cases, the event \textit{Cancel application} never occurs as
the first action of a trace; its most frequent legal position is the
eighth step, occurring there in 287 cases, always after the
application has been processed. The one-step cancellation that
constitutes the entire Pareto front of the unshielded agent is thus
not only structurally illegal under the model but also entirely absent
from the observed behaviour of the process.

\paragraph{Parameter-level generalisation: temporal holdout.}

To evaluate whether the learned policies are specific to a particular
environment parameterisation, the event log was ordered
chronologically by case start timestamp and split, without
shuffling, into a calibration set containing the first 800 cases
(80\%) and a held-out test set containing the final 200 cases
(20\%).

This temporal split follows established practice in process mining
and predictive process monitoring, where random splits can leak
future information into the training partition and therefore
underestimate the impact of process drift.

For each partition, per-activity duration estimates were computed
from the observed executions and used to parameterise an otherwise
identical environment. The process graph, role assignments, and cost
parameters were kept fixed by construction, so that the held-out
environment differs from the calibration environment only in its
empirically observed activity durations.

The six headline $(\alpha,\beta)$ configurations were trained from
scratch on the calibration environment under the standard protocol.
The resulting policies were then frozen and evaluated without any
further training on both the calibration and held-out
parameterisations.

The two parameterisations differ substantially at the activity level.
Table~\ref{tab:duration_drift} reports the per-activity drift between
the calibration and held-out duration estimates, which ranges from
$-11.9\,\%$ to $+13.1\,\%$.

\begin{table}[H]
\centering
\caption{Per-activity duration drift between the calibration (first 800
cases) and held-out test (last 200 cases) parameterisations. Drift is
the relative change in the mean activity duration between the two
partitions.}
\label{tab:duration_drift}
{\renewcommand{\arraystretch}{1.5}
\setlength{\tabcolsep}{6pt}
\begin{tabular}{lrrr}
\toprule
\textbf{Activity} & \textbf{Calib.\ (min)} & \textbf{Test (min)} &
\textbf{Drift} \\
\midrule
Check application form completeness & 54.5  & 54.2  & $-0.5\,\%$ \\
Appraise property                   & 240.0 & 259.6 & $+8.1\,\%$ \\
Check credit history                & 83.2  & 86.2  & $+3.7\,\%$ \\
AML check                           & 251.8 & 284.8 & $+13.1\,\%$ \\
Assess loan risk                    & 178.5 & 163.0 & $-8.7\,\%$ \\
Design loan offer                   & 18.8  & 21.1  & $+12.4\,\%$ \\
Approve loan offer                  & 183.4 & 186.9 & $+1.9\,\%$ \\
Approve application                 & 15.4  & 13.6  & $-11.9\,\%$ \\
Reject application                  & 28.4  & 29.8  & $+4.7\,\%$ \\
Cancel application                  & 13.4  & 14.3  & $+7.1\,\%$ \\
Return application to applicant     & 15.7  & 14.4  & $-8.2\,\%$ \\
Applicant completes form            & 490.5 & 437.7 & $-10.8\,\%$ \\
\bottomrule
\end{tabular}}
\end{table}

The generalisation gap is defined as the difference in the principal
evaluation metrics between the two environments, including mean
episode duration, mean episode cost, Pareto-front structure, and
acceptance rate. A small gap indicates that the policy has learned
stable process behaviour rather than exploiting a particular duration
parameterisation. Table~\ref{tab:generalisation_gap} reports this gap
for each weight configuration.

\begin{table}[H]
\centering
\caption{Temporal generalisation gap under the 80/20 chronological holdout. Policies were trained on the calibration environment (first 800 cases), frozen, and evaluated without retraining on both the calibration and held-out test parameterisations (last 200 cases). Costs are identical by construction because only activity durations differ between environments.}
\label{tab:generalisation_gap}
{\renewcommand{\arraystretch}{1.5}
\setlength{\tabcolsep}{6pt}
\begin{tabular}{lccc}
\toprule
\textbf{Configuration} &
\textbf{Duration (calib., min)} &
\textbf{Duration (test, min)} &
\textbf{Gap} \\
\midrule
$\alpha=0.0,\ \beta=0.0$ & 914.5 & 987.5 & $+8.0\%$ \\
$\alpha=1.0,\ \beta=0.0$ & 1648.0 & 1724.0 & $+4.6\%$ \\
$\alpha=0.0,\ \beta=1.0$ & 412.0 & 431.0 & $+4.6\%$ \\
$\alpha=0.5,\ \beta=0.5$ & 1624.0 & 1703.0 & $+4.9\%$ \\
$\alpha=2.0,\ \beta=0.5$ & 1648.0 & 1724.0 & $+4.6\%$ \\
$\alpha=0.5,\ \beta=2.0$ & 412.0 & 431.0 & $+4.6\%$ \\
\bottomrule
\end{tabular}}
\end{table}
Despite per-activity drifts of up to $\pm 13\,\%$ between the two
parameterisations, the trace-level duration gap remains between
$+4.6\,\%$ and $+8.0\,\%$, because each trace aggregates activities
whose drifts partially cancel. More importantly, the structure of the
recovered front is preserved: pooling the non-baseline configurations
yields the same three non-dominated points in both environments, in the
same dominance ordering, each shifted only along the duration axis, with
no dominance reversal. A policy that had overfitted the calibration
durations would degrade disproportionately under this drift or reorder
the front; neither occurs. The learned trade-off choices therefore
reflect the structure of the process rather than the specific duration
parameterisation used during training, which is the evidence sought
against overfitting. The gap concerns duration only, since costs are
identical by construction between the two environments.

\paragraph{Seed robustness.}

This repeats, under $n=4$ independent random seeds, the same Safe
\ac{RL} experiment whose seed-1 instance is reported as the headline
result of Chapter~\ref{chap:results}
(Table~\ref{tab:loanapp_pareto_summary}). Following standard \ac{RL}
evaluation practice, the full weight sweep was retrained with $n=4$
independent random seeds.
Table~\ref{tab:seed_variance} reports the mean and standard deviation
of the principal metrics across seeds.

\begin{table}[H]
\centering
\caption{Seed robustness of the principal metrics on the Loan
Application benchmark, reported as mean $\pm$ standard deviation across
$n=4$ seeds. Cost is reported under the cost-minimising weight
$(\alpha{=}1,\beta{=}0)$ and duration under the duration-minimising
weight $(\alpha{=}0,\beta{=}1)$, the configurations in which each
objective is the optimisation target and therefore the metric is
informative. The illegal-action ratio is reported twice: pooled across
all six configurations, and excluding the unweighted baseline
$(\alpha{=}0,\beta{=}0)$, which is also excluded from the Pareto
analysis in Chapter~\ref{chap:results}.}
\label{tab:seed_variance}
{\renewcommand{\arraystretch}{1.5}
\setlength{\tabcolsep}{6pt}
\begin{tabular}{lc}
\toprule
\textbf{Metric} & \textbf{Mean $\pm$ std} \\
\midrule
Pareto points recovered            & $26.25 \pm 0.5$ \\
Illegal-action ratio, pooled (\%)  & $1.13 \pm 0.34$ \\
Illegal-action ratio, weighted configs (\%) & $1.15 \pm 0.43$ \\
Mean episode cost $(\alpha{=}1,\beta{=}0)$     & $1503.1 \pm 65.0$ \\
Mean episode duration $(\alpha{=}0,\beta{=}1)$ & $657.3 \pm 48.1$ \\
\bottomrule
\end{tabular}}
\end{table}

\paragraph{Pareto front stability.}

The recovered Pareto front is robust to the random seed in both range
and shape. The extreme points of the front (minimum cost and minimum
duration) are identical across all four seeds, and the intermediate
points coincide almost point-for-point. This contrasts with the
Shielding-only baseline analysed in Section~\ref{sec:shield_contribution},
where random valid-action sampling recovered the same broad trade-off shape but showed seed-dependent variation in the number of Pareto points and in the front extremes.

Figure~\ref{fig:pareto_overlay} overlays the four fronts; they are so
close that each seed occludes the ones plotted before it, leaving only a single visible curve except for a narrow mid-cost region around
$1900$--$2050$, where two seeds trace marginally different paths. The
legend's per-seed point count (26 for three seeds, 27 for the fourth)
is the same small variation already reported as $26.25 \pm 0.5$ in
Table~\ref{tab:seed_variance}.


\paragraph{Compliance stability across seeds.}

Across the four seeds, the final illegal-action ratio is low and
consistent for every weight configuration: pooled across all six
configurations it averages $1.67 \pm 0.89\,\%$, and restricted to the
five weighted configurations it averages $1.88 \pm 1.07\,\%$. Two
single-seed values stand out from an otherwise tight distribution
(seed 3 at $\alpha{=}0.5,\beta{=}2.0$: $9.15\,\%$; seed 4 at
$\alpha{=}2.0,\beta{=}0.5$: $11.58\,\%$), but both remain well below
the unweighted baseline's per-seed range of $0.41$--$0.85\,\%$, and
every seed still reaches $100\,\%$ acceptance in the final $20\,\%$ of
training, with no episode left incomplete by truncation
(Figures~\ref{fig:illegal_boxplot} and~\ref{fig:acceptance_boxplot}).
The unweighted baseline $(\alpha{=}0,\beta{=}0)$ behaves no differently
from the weighted configurations overall.

\begin{figure}[H]
    \centering
    \includegraphics[width=0.80\linewidth]{Pictures/seed_robustness_pareto_overlay.png}
    \caption{Pareto fronts recovered by Safe \ac{RL} under $n=4$ random
seeds on the Loan Application benchmark, with the recovered point
count per seed given in the legend. The fronts are plotted
in seed order and are near-coincident in both range and
shape, so later seeds occlude earlier ones; the only
visible divergence is a narrow mid-cost region around
$1900$--$2050$. Compared with the Shielding-only seed
analysis in Section~\ref{sec:shield_contribution}, this
near-perfect overlap indicates that the learned front is
not merely a product of stochastic valid-trace sampling, but
a stable outcome of the shielded learning framework.}
    \label{fig:pareto_overlay}
\end{figure}

\begin{figure}[H]
    \centering
    \includegraphics[width=0.82\linewidth]{Pictures/seed_robustness_illegal_boxplot.png}
    \caption{Distribution of the final illegal-action ratio per $(\alpha,\beta)$ configuration across $n=4$ seeds, combining boxplots with explicit seed tracking. Individual runs are color-coded to isolate variance. Every configuration and seed clusters below $8\,\%$.}
    \label{fig:illegal_boxplot}
\end{figure}

\begin{figure}[H]
    \centering
    \includegraphics[width=0.82\linewidth]{Pictures/seed_robustness_acceptance_boxplot.png}
    \caption{Total number of accepting episodes recorded within the final 20\,\% of training per configuration across $n=4$ seeds. Every run, including the unweighted baseline, reaches a stable plateau of approximately $2\,700$--$2\,930$ accepting episodes.}
    \label{fig:acceptance_boxplot}
\end{figure}



\paragraph{Scope of the generalisation claim.}
Two qualifications bound the claim supported by these results.
First, robustness is established across random seeds on a
fixed process graph; transfer to unseen process
configurations is not evaluated here, since the shield is
synthesised from the same \ac{DCR} model used for training and
the policies are therefore process-specific by design. This
particular question, whether the framework's behaviour generalises to
process structures never seen during development, is instead addressed
independently in Chapter~\ref{chap:external_validation}. Second, this
seed-robustness experiment was run on the Loan Application benchmark
only; the Expense Report benchmark was evaluated under a single seed
($s=1$) and was not repeated across seeds. The per-event role-assignment
stability of the role-conditioned experiment, by contrast, was measured
directly across the same four seeds, not merely inferred from the
front-level statistics above (Section~\ref{sec:discussion_roles}).
The implications of both qualifications are discussed in
Section~\ref{sec:limitations}.
